% TOP_PROBLEM: {"id": 3000030, "title": "Disjoint spanning in- and out-arborescences", "category_id": 2, "category": {"id": 2, "name": "combinatorics", "display_name": "Combinatorics", "description": "Counting problems, graph theory, discrete structures.", "slug": "combinatorics", "order_index": 2, "created_at": "2026-07-31T15:26:25.671Z"}, "status": "open", "problem_number": "AMR-029-0030", "set_id": 13}
% FIRST_POSTED: 2026-10-01
% ENTRY_KIND: statement_only
% UnsolvedMath ID 3000030; source catalogue code AMR-029-0030
% !TeX program = lualatex
% Definition and sources only; this file is not an LLM solution attempt.
\documentclass[11pt,a4paper]{article}
\usepackage[margin=25mm]{geometry}
\usepackage{fontspec}
\setmainfont{lmroman10-regular.otf}[BoldFont=lmroman10-bold.otf,
 ItalicFont=lmroman10-italic.otf,BoldItalicFont=lmroman10-bolditalic.otf]
\usepackage{amsmath,amssymb,mathtools}
\usepackage{unicode-math}
\setmathfont{latinmodern-math.otf}
\AtBeginDocument{\let\setminus\smallsetminus}
\usepackage{xurl,hyperref}
\hypersetup{colorlinks=true,urlcolor=blue}
\setcounter{secnumdepth}{0}
\setlength{\parindent}{0pt}
\setlength{\parskip}{5pt}
\setlength{\emergencystretch}{3em}
\begin{document}
\section*{Disjoint spanning in- and out-arborescences}\label{3000030}
{\small UnsolvedMath ID 3000030 · AMR-029-0030 · Combinatorics · AMR Open Problem Lists}
\subsection{Definitions and mathematical statement}
A finite loopless digraph $D=(V,A)$ is $k$-arc-connected if every
nonempty proper subset $X\subset V$ has at least $k$ entering arcs.
Applying the same condition to $V\setminus X$ also gives at least
$k$ leaving arcs. Equivalently, every ordered pair of distinct vertices
can be joined by $k$ pairwise arc-disjoint directed paths.

For a chosen root $r$, a spanning out-arborescence is a directed spanning
tree with a directed path from $r$ to every vertex. A spanning
in-arborescence has a directed path from every vertex to $r$.
Is there an absolute positive integer $k_0$ such that every
$k_0$-arc-connected digraph, for every choice of $r\in V$, contains
an in-arborescence $T^-$ and an out-arborescence $T^+$ with
\[
                             A(T^-)\cap A(T^+)=\varnothing?
\]
Both trees must span the same full vertex set and use the same prescribed
root. They may share all their vertices; disjointness concerns arcs.
The required constant cannot depend on the size or structure of the
digraph, or on the selected root. Separate existence of an incoming
tree and an outgoing tree follows under much weaker connectivity,
but does not establish that both can be chosen simultaneously without
sharing an arc. Parallel arcs, if present, are separate available arcs.
\subsection{Short English statement}
Does sufficiently high arc connectivity always guarantee disjoint
incoming and outgoing spanning trees at every prescribed root?
\subsection{Sources}
\begin{itemize}
\item[S1] ULAM.AI and the UnsolvedMath Contributors, supplied problems.json, record 3000030 (AMR-029-0030), AMR Open Problem Lists. Catalogue identity and source transcription; CC BY 4.0.
\url{https://huggingface.co/datasets/ulamai/UnsolvedMath}.
\item[S2] Egres Open - List of open problems, Open-problem page: Disjoint spanning in- and out-arborescences. Original source indexed by AMR-029-0030.
\url{https://oldlemon.cs.elte.hu/egres/open/List_of_open_problems}.
\item[S3] EGRES Open, Disjoint spanning in- and out-arborescences. Individual source page and explanatory remarks.
\url{https://lemon.cs.elte.hu/egres/open/Disjoint_spanning_in-_and_out-arborescences}.
\end{itemize}
\end{document}
